\section{Thesis Outline}\label{sec:intro_outline}

The remainder of this thesis is structured as follows.
Chapter~\ref{ch:fundamentals} introduces the background required to follow
the rest of the work: the foundations of knowledge graphs and graph-based
machine learning, explainable AI, and the BMW quality management and 8D
problem-solving process in which this thesis is situated.
Chapter~\ref{ch:related_work} reviews existing approaches to causer and
root-cause attribution -- rule-based systems, black-box machine learning,
and knowledge-graph retrieval -- and derives the research gap that
motivates this thesis.
Chapter~\ref{ch:problem} formalizes causer attribution as an inductive
link-prediction problem on an event-centric knowledge graph and describes
the underlying data sources and graph schema.
Chapter~\ref{ch:method} presents the proposed method: the model
architectures compared in this work (rule-based baseline, knowledge graph
embeddings, and a relational graph neural network) and the evidence
retrieval and explanation generation pipeline built on top of them.
Chapter~\ref{ch:setup} details the experimental setup, including the
baselines, the inductive train/test protocol with leakage controls, and
the evaluation metrics used to assess predictive performance and
explanation faithfulness.
Chapter~\ref{ch:results} reports the quantitative results, a qualitative
analysis of the generated explanations, and the limitations of the
proposed approach.
Chapter~\ref{ch:conclusion} concludes the thesis by revisiting the research
questions and contributions in light of the results and outlining
directions for future work.
